\subsubsection{The eo\_pipe Software Package}

The eo\_pipe software package (\url{https://github.com/lsst-camera-dh/eo_pipe}) is used to perform a standard set of electro-optical tests on LSSTCam pixel data. The initial set of tests were originally implemented for acceptance testing of the CCDs provided by the ITL and e2V vendors.  The criteria and the algorithms used to analyze the CCD acceptance data are described in LCA-128 (\url{https://ls.st/LCA-128}) and LCA-10103 (\url{https://ls.st/LCA-10103}), respectively.  The acceptance measurements included read noise, non-linearity, serial and parallel CTE, crosstalk, dark current, quantum efficiency, bright and dark pixel defects, traps, pixel response non-uniformity, and point spread function sizes.  Over time, as testing progressed from CCD to raft to focal plane level, various other measurements and algorithms were added: photon transfer curve (PTC) generation, tearing detection, divisadero tearing, brighter-fatter analysis, bias shift detection, bias stability, CTI as a function of flux, gain stability, persistence measurement, row-means correlations, raft-level overscan and imaging region correlations, and scan mode analysis.

To ensure consistency with the official calibration products, the cp\_pipe pipelines (\url{https://github.com/lsst/cp_pipe}) are run first to generate various data products that are inputs to the eo\_pipe pipelines, i.e., combined bias, dark, and flat images, catalogs of pixel defects, PTCs, etc..  The eo\_pipe tasks also perform other processing of the raw images to fill out its suite of tests.  The eo\_pipe analyses are implemented as "pipeline tasks" using the Rubin Science Pipelines middleware.  This allows them to be run on LSSTCam data that has been ingested into a Rubin data repository and accessed by the Rubin data butler.   The Rubin Batch Production Service framework (\url{https://github.com/lsst/ctrl_bps}) is used to run the eo\_pipe pipelines at scale at the US Data Facility using the slurm batch system.

The eo\_pipe results are persisted into the same data repository containing the raw LSSTCam data in collections that are labeled by analysis type, "acquisition run", and weekly software release.  Subsequent specialized follow up analyses can then be run on the eo\_pipe data products by accessing those collections.  The acquisition runs are identified for each set of LSSTCam configurations, comprising sequencer versions and various other LSSTCam conditions and settings.  The data from those acquisition runs are processed together in an "analysis run" using specific versions of the cp\_pipe and eo\_pipe pipelines.  Standard outputs from these analysis runs are available in web-based reports at \url{https://s3df.slac.stanford.edu/data/rubin/lsstcam/}